\documentclass{article}
\usepackage[utf8]{inputenc}
\usepackage{booktabs}
\usepackage{longtable}
\title{AAP live one-paper}
\begin{document}
\maketitle
\begin{abstract}
AAP live one-paper synthesizes 1 papers across 5 evidence rows using factor structure as the main organizing lens. The synthesis is concentrated in Survey data from 260 respondents in Dalian, China, behavioral factors, with recurring methods such as Structural Equation Modeling (SEM), Theory of Planned Behavior (TPB) extension and task coverage centered on Predict pedestrians' violating road-crossing behavior intentions. Across 1 papers, the evidence is most coherently organized around factor structure rather than any single method, and the strongest clusters are Survey data from 260 respondents in Dalian, China, behavioral factors, demographic factors. Future work should prioritize standardized reporting and more explicit cross-study comparison.
\end{abstract}
\noindent\textbf{Keywords:} structural equation modeling (sem), theory of planned behavior (tpb) extension, predict pedestrians' violating road-crossing behavior intentions, survey data from 260 respondents in dalian, china, examine pedestrians' self-reported violating crossing behavior intentions by applying and extending the theory of planned behavior (tpb), testing instrumental attitude, subjective norm, perceived behavioral control, and extended factors including descriptive norm, perceived risk, and conformity tendency, cronbach's alpha (reliability), confirmatory factor analysis (validity), cmin/df ratio, behavioral factors
\section{Introduction and Research Context}
Pedestrian-road traffic collisions remain a leading cause of unintentional injury mortality globally, with intersections representing hotspots where conflicting pedestrian and vehicle trajectories elevate crash risk. Among the behavioral factors contributing to pedestrian injuries at signalized intersections, violating crossing behavior—defined as crossing against pedestrian traffic signals or outside designated crossing areas—stands out as a prevalent and preventable risk factor. Empirical traffic observation studies consistently document substantial proportions of pedestrians engaging in red-light violations, particularly in high-density urban environments where waiting delays compete with time constraints and perceived convenience. The consequences are severe: pedestrian violations account for a disproportionate share of intersection-related casualties, with injury severity amplified by the physical vulnerability of pedestrians relative to vehicle occupants. \cite{08_2016_zhou_aap}

Understanding the determinants of pedestrian violating crossing behavior has attracted growing research attention from transportation safety scholars, yet the literature has historically lacked unified theoretical frameworks capable of integrating the multiple cognitive, social, and contextual influences on crossing decisions. Early approaches tended to treat crossing behavior as a function of situational factors such as waiting time, vehicle gap size, and signal timing, while underweighting the psychological mechanisms through which pedestrians evaluate social norms, personal attitudes, and perceived control over their actions. This fragmentation has limited the translational utility of research findings for designing effective interventions. \cite{08_2016_zhou_aap}

This review synthesizes evidence from a systematic examination of the Theory of Planned Behavior (TPB) as applied to pedestrian violating crossing intentions, drawing on a primary study conducted in Dalian, China, that employed structural equation modeling to test both basic and extended TPB specifications. The review's scope encompasses behavioral, demographic, and traffic-related predictors of crossing intentions, with particular attention to how TPB constructs perform in predicting self-reported violating behavior and which extensions improve model explanatory power. The contribution lies in characterizing the relative importance of instrumental attitudes, descriptive and subjective norms, perceived behavioral control, conformity tendencies, and perceived risk in shaping pedestrians' intentions to cross against traffic signals, thereby identifying leverage points for safety interventions. \cite{08_2016_zhou_aap}

\section{Theoretical Background: Theory of Planned Behavior and Extensions}
The Theory of Planned Behavior, developed by Ajzen (1991), provides a cognitive-behavioral framework for predicting intentional behaviors across diverse domains. The theory posits that behavioral intention—the immediate antecedent of actual behavior—is a function of three primary constructs: attitude toward the behavior, subjective norm, and perceived behavioral control. Attitude reflects the individual's evaluation of the behavior as favorable or unfavorable, encompassing both instrumental dimensions (e.g., beneficial versus harmful) and experiential dimensions (e.g., enjoyable versus unpleasant). Subjective norm captures the perceived social pressure to perform or not perform the behavior, derived from perceptions of what important referent others think one should do. Perceived behavioral control reflects the individual's belief about their ability to execute the behavior, analogous to Bandura's concept of self-efficacy insofar as it concerns perceived ease or difficulty of performance. \cite{08_2016_zhou_aap}

The theory's applicability to health-related and risky behaviors has been extensively validated, but its application to pedestrian safety represents a more recent development. Pedestrian violating crossing behavior presents a distinctive case because the behavior is socially discouraged yet frequently enacted, involves competing considerations of time efficiency and safety, and is subject to social observation that may trigger normative influences beyond traditional subjective norms. These characteristics have motivated researchers to extend the basic TPB model with additional constructs that capture dimensions of pedestrian decision-making not fully represented in the original framework. \cite{08_2016_zhou_aap}

Three extensions have received particular attention in the pedestrian behavior literature. First, descriptive norm has been proposed as a complement to subjective norm, capturing perceptions of what others actually do rather than what others think one should do. The distinction is meaningful because pedestrians may observe others violating traffic signals and infer that such behavior is common, even while believing that referents would disapprove of violations personally. Second, conformity tendency has been introduced as an individual difference variable reflecting the propensity to align one's behavior with observed social behavior, particularly in ambiguous or uncertain situations. Third, perceived risk has been incorporated to represent the cognitive appraisal of traffic danger, including perceived susceptibility to injury and perceived severity of potential consequences. The study under review, Zhou et al. (2016), systematically tested each of these extensions alongside the basic TPB constructs to evaluate their incremental predictive utility for pedestrian violating crossing intentions. \cite{08_2016_zhou_aap}

\section{Structural Equation Modeling Methodology}
The methodological approach employed in the reviewed study reflects a growing trend in behavioral traffic research toward structural equation modeling (SEM) as a technique for testing complex, multi-path theoretical models. SEM enables simultaneous estimation of multiple linear relationships between latent constructs (theoretical variables) and their observed indicators (survey items), thereby allowing researchers to evaluate whether the hypothesized directional relationships among constructs are supported by the data while accounting for measurement error. This represents an advantage over traditional regression-based approaches that estimate individual paths in isolation and do not formally incorporate the measurement model. \cite{08_2016_zhou_aap}

The study by Zhou et al. (2016) implemented SEM using questionnaire survey data collected from 260 respondents in Dalian, China. Participants responded to a scenario-based instrument presenting a situation in which crossing against pedestrian traffic lights was possible, then indicated their intentions, attitudes, norms, and control perceptions regarding such behavior. The analytical procedure proceeded in stages: first, data reliability was assessed using Cronbach's alpha coefficients to verify internal consistency across multi-item scales; second, confirmatory factor analysis was conducted to establish the validity of the measurement model by verifying that observed items loaded onto their intended latent factors; third, structural models were estimated to test the hypothesized relationships between latent predictors and the behavioral intention outcome. \cite{08_2016_zhou_aap}

Four nested models were tested to enable model comparison: the basic TPB specification containing instrumental attitude, subjective norm, and perceived behavioral control as predictors; the basic TPB plus descriptive norm; the basic TPB plus perceived risk; and a full extended model containing instrumental attitude, descriptive norm, and conformity tendency. Goodness of fit was evaluated using multiple indices, including the CMIN/df ratio, the Root Mean Square Error of Approximation (RMSEA), and the Comparative Fit Index (CFI). The final model demonstrated excellent fit with a CFI of 0.996, indicating that the extended TPB specification effectively captured the data structure. This methodological approach provides a rigorous test of theoretical predictions while acknowledging the latent-variable nature of psychological constructs, though the reliance on self-reported intention as the outcome variable introduces limitations that merit acknowledgment. \cite{08_2016_zhou_aap}

\section{Behavioral Factors in Crossing Intentions}
The reviewed evidence reveals substantial heterogeneity in the predictive power of different behavioral factors, with some TPB constructs performing robustly while others fail to reach statistical significance. Instrumental attitude emerged as the most consistent predictor of pedestrian violating crossing intention across all four models tested, indicating that the degree to which pedestrians evaluate crossing against traffic signals as beneficial, effective, or advantageous carries significant weight in their behavioral intentions. This finding aligns with the broader TPB literature suggesting that attitude typically constitutes the strongest direct predictor of intention when the behavior is under volitional control. \cite{08_2016_zhou_aap}

A notable methodological finding concerns the relative predictive strength of descriptive norm compared to subjective norm. The study found that descriptive norm was a stronger predictor than subjective norm, and critically, when descriptive norm was added to the model, subjective norm became non-significant. This pattern suggests that what pedestrians believe others actually do (the descriptive norm) carries more behavioral weight than what pedestrians believe others think they should do (the subjective norm), at least in the context of traffic signal violations. Given that 93.1 percent of respondents disapproved of violating behavior (family mean = 2.93, friends mean = 3.20 on a 7-point scale), the disconnect between personal disapproval and observed behavior among others appears to drive the predictive power of descriptive norm. \cite{08_2016_zhou_aap}

Conformity tendency was a significant predictor in the extended model (Model 4), indicating that pedestrians higher in the tendency to conform to others' behavior are more likely to express intentions to cross against traffic signals when they observe others doing so. This finding points to a social learning mechanism in which violating behavior propagates through observational exposure, with individual differences in conformity amplifying the effect. By contrast, perceived risk was not a significant predictor despite respondents reporting high perceived risk of injury (mean = 5.77) and danger to life (mean = 5.83) on 7-point scales. The authors interpreted this non-significant finding as indicating that pedestrians fail to adequately perceive the susceptibility and seriousness of risks associated with traffic violations—the high risk ratings may reflect abstract acknowledgment of danger without translation into behavioral guidance. Perceived behavioral control was also non-significant across all models, likely because crossing the street is not perceived as challenging for able-bodied pedestrians in typical urban environments, rendering this construct moot for predicting violations. \cite{08_2016_zhou_aap}

\section{Demographic and Traffic Context Factors}
While the primary focus of the reviewed study lies in testing TPB constructs, the sample characteristics provide insight into demographic patterns that may shape crossing intentions. The sample of 260 respondents in Dalian, China, comprised 57.3 percent female and 42.7 percent male participants, with age distribution concentrated in younger cohorts: 49.2 percent were aged 18-24 years and 33.1 percent were aged 25-39 years. This demographic profile—predominantly young and slightly female-skewed—limits the ability to examine age or gender differences in violating crossing intentions within the study itself, as the sample does not include sufficient representation of older pedestrians or balance across genders to enable meaningful subgroup analysis. \cite{08_2016_zhou_aap}

The traffic context in Dalian represents a high-density urban Chinese intersection environment where pedestrian volumes are substantial and signalized crossing opportunities are frequent. The scenario-based methodology presented participants with a hypothetical situation involving crossing against pedestrian traffic lights, asking them to report their likely intentions under described conditions. While this approach enables standardized measurement of psychological constructs, it necessarily abstracts from the real-time dynamics of actual crossing decisions, where situational factors such as vehicle speed, gap size, waiting time, and presence of other pedestrians fluctuate continuously. The study did not directly incorporate traffic flow variables or intersection characteristics into the structural model, instead focusing on individual-level psychological predictors. \cite{08_2016_zhou_aap}

The concentration of the sample in a single Chinese city raises generalizability considerations for the findings. Pedestrian crossing behavior in China may differ from behavior in other cultural contexts due to differences in traffic density, enforcement practices, cultural attitudes toward rule compliance, infrastructure design, and the prevalence of electric bicycle and e-bike traffic that creates distinct crossing environments. The authors explicitly noted this limitation, acknowledging that findings may not generalize to other cultural contexts where pedestrian crossing behavior differs. Future research employing samples from diverse geographic and cultural settings would help establish the external validity of the TPB extensions identified as significant predictors. \cite{08_2016_zhou_aap}

\section{Study Limitations and Methodological Considerations}
The reviewed research is subject to several methodological limitations that constrain the interpretability and generalizability of its findings. Most fundamentally, the reliance on self-reported behavioral intention as the dependent variable introduces a common-method bias concern, as both predictors and outcome were measured from the same respondents at the same time point using the same questionnaire instrument. While the structural model tests directional relationships between latent constructs, the correlational nature of cross-sectional survey data precludes strong causal inference. Actual behavior may diverge from stated intention due to unobserved situational constraints, social desirability pressures, or changes in contextual factors that alter the decision context between intention formation and behavior execution. \cite{08_2016_zhou_aap}

The relatively small number of items in the questionnaire, maintained to ensure internal reliability of each scale, may have constrained the content validity of the measurement instrument. Each latent construct was operationalized using a limited number of indicator items, potentially missing important facets of the underlying psychological dimension. The sample size of 260, while adequate for SEM estimation with the number of parameters estimated, limits statistical power for detecting small effect sizes and precludes meaningful multi-group analyses across demographic subgroups. \cite{08_2016_zhou_aap}

The cultural specificity of the sample—drawn exclusively from Dalian, China—represents perhaps the most consequential limitation for the field's cumulative knowledge. Pedestrian violating crossing behavior is known to vary substantially across cultural contexts, with some environments exhibiting higher prevalence of red-light violations due to different norm enforcement regimes, traffic densities, and infrastructure designs. The finding that descriptive norm outperformed subjective norm may be particularly context-dependent, reflecting the observed behavior of other pedestrians in Chinese urban traffic environments where violations may be more visible or common. Whether these patterns replicate in European, North American, or other Asian contexts remains an open empirical question that the current evidence base cannot address. \cite{08_2016_zhou_aap}

\section{Conclusions and Implications for Safety Interventions}
The reviewed evidence from Zhou et al. (2016) demonstrates that the Theory of Planned Behavior, when extended with descriptive norm and conformity tendency, provides a credible framework for predicting pedestrian violating crossing intentions. Instrumental attitude emerged as the most robust predictor across model specifications, indicating that pedestrians' evaluations of the personal benefits and costs of violating behavior substantially shape their intentions. The finding that descriptive norm outperforms subjective norm in predictive power is particularly noteworthy: it suggests that safety interventions targeting observed social behavior may be more effective than those emphasizing injunctive norms (what others think one should do). If pedestrians calibrate their behavior to what they see others doing rather than what they believe others expect of them, then reducing the visibility of violations through enforcement or infrastructure design could produce spillover effects on other pedestrians' behavior. \cite{08_2016_zhou_aap}

The non-significance of perceived risk despite high reported risk levels points to a critical disconnect in pedestrian cognition—acknowledgment of danger does not translate into risk-averse decision-making under the time pressure and convenience considerations that characterize real-world crossing decisions. This finding implies that safety campaigns emphasizing danger severity may have limited impact if they do not address the perceived benefits of violation or the social modeling of violations in the environment. The non-significance of perceived behavioral control further suggests that interventions targeting self-efficacy or control beliefs are unlikely to yield gains in this behavioral domain, as crossing is not subjectively perceived as difficult for most pedestrians. \cite{08_2016_zhou_aap}

For traffic safety interventions, these findings suggest that structural changes to the traffic environment—such as reducing waiting times, improving signal timing, or designing infrastructure that minimizes the visibility of others' violations—may be more effective than individual-level psychological interventions. However, the cultural specificity of the evidence base requires caution in generalization. The field would benefit from replication studies in diverse cultural and traffic environments, longitudinal designs that track the translation of intention to actual behavior, and experimental tests of intervention effectiveness grounded in the identified predictors. \cite{08_2016_zhou_aap}

\appendix
\section{Appendix}
\subsection{Appendix Summary}
\noindent The appendix records 1 papers and 5 matrix rows used to generate the review synthesis. The most visible methodological families are Structural Equation Modeling (SEM), Theory of Planned Behavior (TPB) extension, while recurring tasks include Predict pedestrians' violating road-crossing behavior intentions.

\begin{itemize}
\item \textbf{Papers}: 1
\item \textbf{Matrix rows}: 5
\item \textbf{Verified gaps}: 
\item \textbf{Dominant axis}: factor
\item \textbf{Top methods}: Structural Equation Modeling (SEM), Theory of Planned Behavior (TPB) extension
\item \textbf{Top tasks}: Predict pedestrians' violating road-crossing behavior intentions
\item \textbf{Top datasets}: Survey data from 260 respondents in Dalian, China
\end{itemize}

\subsection{Evidence Table}
\begin{longtable}{p{0.14\textwidth}p{0.21\textwidth}p{0.10\textwidth}p{0.17\textwidth}p{0.17\textwidth}p{0.17\textwidth}}
\toprule
Paper & Title & Year & Methods & Tasks & Gap matches \\
\midrule
\endfirsthead
\toprule
Paper & Title & Year & Methods & Tasks & Gap matches \\
\midrule
\endhead
08\_2016\_zhou\_aap & An extension of the theory of planned behavior to predict pedestrians' violating crossing behavior using structural equation modeling & 2016 & Structural Equation Modeling (SEM), Theory of Planned Behavior (TPB) extension & Predict pedestrians' violating road-crossing behavior intentions & -- \\
\bottomrule
\end{longtable}
\end{document}
% BIB START
@article{08_2016_zhou_aap,
  author={Hongmei Zhou; Stephanie Ballon Romero; Xiao Qin},
  title={An extension of the theory of planned behavior to predict pedestrians' violating crossing behavior using structural equation modeling},
  year={2016},
  journal={Accident Analysis and Prevention}
}